% Paper 6 regime diagram: exact g(rho,2) vs falsified g-tilde (Panel A) and
% the succession price D* (Panel B), recomputed directly from
% create_regime_diagram() in figures/src/paper6_figures.py. Curves sampled
% at reduced density (25 pts Panel A, 20 pts Panel B) from the same closed
% forms the script evaluates at full density (300 pts); the sample grid is
% chosen to land exactly on the crossing point (rho=(3-sqrt5)/2 in A,
% D*=0.34962 in B) so the shaded-region polygons meet the curves cleanly.
% Deterministic closed-form; the script's np.random.seed(20260816) is set
% but unused by either panel (no sampling here).
% NOTE (flagged, not silently patched): the script defines erlang_c(c,rho)
% but create_regime_diagram() never calls it -- Panel B's K=1.0362 is a
% hardcoded literal (matching the paper's own by-hand computation in
% Sec. "Numbers by hand", confirmed independently below), not derived
% programmatically from erlang_c(). Separately, erlang_c() itself is dead
% code under the current numpy (np.math.factorial no longer exists --
% numpy.math was removed from the numpy release pinned in this container),
% so if anything ever calls it as written, it raises AttributeError. Since
% it is not on the current execution path for either PNG or this figure,
% it does not affect the plotted data -- but it is worth a fix upstream.
% Independent check performed for this conversion: erlang_c(4,0.20833)
% (via math.factorial, sidestepping the numpy bug) = 0.011024, matching the
% paper's cited C(4,0.208)=0.0110 and K=0.2+0.8362=1.0362 [verified,
% recomputed for this conversion]. dstar=eta*K/(1-eta*K)=0.34962 (paper:
% D*=0.350, rounded); winf(1.5*dstar)=1.37607 (paper: 1.376). All match.
% [verified: script paper6_figures.py, deterministic closed-form]
\begin{figure}[t]
\centering
\begin{minipage}{0.48\textwidth}
\centering
\begin{tikzpicture}
\begin{axis}[
  harbor curve,
  xlabel={utilization $\rho$},
  ylabel={skill premium $r=\mu_s/\mu_g$},
  title={\textbf{A} --- the exact boundary\\vs.\ the falsified one},
  xmin=0,xmax=0.95,ymin=0.9,ymax=4.0,
  xtick={0,0.2,0.4,0.6,0.8},
  ytick={1.0,1.5,2.0,2.5,3.0,3.5,4.0},
  legend pos=north west,
]
% The two lenses were separated by hue alone (shipred vs harborblue fill at
% ~0.15 opacity), which is the least accurate visual channel and collapses to one
% grey in print -- and neither carried a text label, so the caption was doing the
% figure's work. Both now carry an opposed hatch AND an in-region text label, so
% "these are two different kinds of error" survives the delete-all-colour test.
\addplot[draw=none,fill=white,pattern=north east lines,pattern color=shipred,forget plot] coordinates {
  (0.0200,1.0396) (0.0562,1.1092) (0.0924,1.1762) (0.1286,1.2406) (0.1648,1.3024)
  (0.2010,1.3616) (0.2372,1.4181) (0.2734,1.4720) (0.3096,1.5233) (0.3458,1.5720)
  (0.3820,1.6180)
  (0.3820,1.6180) (0.3458,1.5285) (0.3096,1.4484) (0.2734,1.3762) (0.2372,1.3109)
  (0.2010,1.2515) (0.1648,1.1973) (0.1286,1.1476) (0.0924,1.1018) (0.0562,1.0595)
  (0.0200,1.0204)
} \closedcycle;
\addplot[draw=none,fill=white,pattern=north west lines,pattern color=harborblue,forget plot] coordinates {
  (0.3820,1.6180) (0.4204,1.7253) (0.4588,1.8478) (0.4973,1.9891) (0.5357,2.1537)
  (0.5741,2.3481) (0.6126,2.5810) (0.6510,2.8652) (0.6894,3.2197) (0.7278,3.6744)
  (0.7663,4.0000) (0.8047,4.0000) (0.8431,4.0000) (0.8816,4.0000) (0.9200,4.0000)
  (0.9200,1.9936) (0.8816,1.9860) (0.8431,1.9754) (0.8047,1.9619) (0.7663,1.9454)
  (0.7278,1.9259) (0.6894,1.9035) (0.6510,1.8782) (0.6126,1.8499) (0.5741,1.8186)
  (0.5357,1.7844) (0.4973,1.7473) (0.4588,1.7071) (0.4204,1.6641) (0.3820,1.6180)
} \closedcycle;
\addplot[seagreen,thick] coordinates {
  (0.0200,1.0396) (0.0562,1.1092) (0.0924,1.1762) (0.1286,1.2406) (0.1648,1.3024)
  (0.2010,1.3616) (0.2372,1.4181) (0.2734,1.4720) (0.3096,1.5233) (0.3458,1.5720)
  (0.3820,1.6180) (0.4204,1.6641) (0.4588,1.7071) (0.4973,1.7473) (0.5357,1.7844)
  (0.5741,1.8186) (0.6126,1.8499) (0.6510,1.8782) (0.6894,1.9035) (0.7278,1.9259)
  (0.7663,1.9454) (0.8047,1.9619) (0.8431,1.9754) (0.8816,1.9860) (0.9200,1.9936)
};
\addlegendentry{exact $g(\rho,2)=1+2\rho-\rho^2$}
\addplot[shipred,thick,dashed] coordinates {
  (0.0200,1.0204) (0.0562,1.0595) (0.0924,1.1018) (0.1286,1.1476) (0.1648,1.1973)
  (0.2010,1.2515) (0.2372,1.3109) (0.2734,1.3762) (0.3096,1.4484) (0.3458,1.5285)
  (0.3820,1.6180) (0.4204,1.7253) (0.4588,1.8478) (0.4973,1.9891) (0.5357,2.1537)
  (0.5741,2.3481) (0.6126,2.5810) (0.6510,2.8652) (0.6894,3.2197) (0.7278,3.6744)
  (0.7663,4.2786) (0.8047,5.1205) (0.8431,6.3750) (0.8816,8.4437) (0.9200,12.5000)
};
\addlegendentry{proposed $\tilde g=1+\rho/(1-\rho)$ (falsified)}
\addplot[black,dotted,forget plot] coordinates {(0.381966,0.9) (0.381966,4.0)};
% Moved off the legend: legend pos=north west with these two long
% (math-bearing) entries renders nearly full plot width and occupies
% roughly the top fifth of the y-range, which is exactly where this label
% sat before (y=3.62, just under ymax=4.0) -- the crossing marker's own
% column, not something the shared axis style touches. Lower + smaller
% keeps it clear of both the legend above and the false-certify cluster
% (~y=1.0-1.15) below.
\node[font=\tiny,rotate=90,anchor=north] at (axis cs:0.394,2.3) {$\rho=(3-\sqrt5)/2$};
% Truncation flagged rather than silent: the dashed curve leaves the top of the
% frame at rho~0.755 and reaches 12.5 by rho=0.92. Panel A cannot show the
% saturate-vs-diverge contrast (that is a c=8, log-axis picture); it can at least
% not pretend the dashed line stopped.
\node[shipred,font=\tiny,align=right,anchor=north east,fill=white,fill opacity=0.88,text opacity=1,inner sep=1.5pt]
  at (axis cs:0.945,3.96) {axis clipped at $r{=}4$;\\$\tilde g$ reaches $12.5$ at $\rho{=}0.92$};
% Region labels: each lens is named ON the figure, not only in the caption.
\addplot[shipred,mark=x,mark size=3.0pt,mark options={line width=1.1pt},only marks,forget plot] coordinates {(0.1,1.15)};
\node[shipred,font=\tiny,align=center,anchor=west,fill=white,fill opacity=0.88,text opacity=1,inner sep=1.5pt]
  at (axis cs:0.135,1.055) {\textbf{false-certify} lens\\$\tilde g$ certifies $r{=}1.15$, $g$ rejects};
\draw[shipred,->,thin] (axis cs:0.24,1.20) -- (axis cs:0.24,1.31);
\addplot[harborblue,mark=o,mark size=2.4pt,mark options={line width=1.1pt},only marks,forget plot] coordinates {(0.5,1.9)};
\node[harborblue,font=\tiny,align=center,anchor=south east,fill=white,fill opacity=0.88,text opacity=1,inner sep=1.5pt]
  at (axis cs:0.90,2.05) {\textbf{false-reject} lens\\$\tilde g$ rejects $r{=}1.9$, $g$ accepts};
\end{axis}
\end{tikzpicture}
\end{minipage}%
\hfill
\begin{minipage}{0.48\textwidth}
\centering
\begin{tikzpicture}
\begin{axis}[
  harbor curve,
  xlabel={death-rate $\times$ succession-time, $\xi/\eta$},
  % Neither plotted series is a response time: one is an infimum OVER response
  % times, the other a cost line that only carries time units because A/(w*lambda)
  % does. The old ylabel invited reading the seagreen curve as the specialist's
  % actual W -- the one misreading this panel most needs to prevent.
  ylabel={response time / cost line (same units)},
  title={\textbf{B} --- the succession price\\($\eta{=}0.25$, $K{=}1.036$)},
  % ymax raised from 2.6 to 4.2 so the finite-premium W_bd series fits in frame.
  xmin=0,xmax=1.0,ymin=0,ymax=4.2,
  xtick={0,0.2,0.4,0.6,0.8,1.0},
  ytick={0,1,2,3,4},
  legend pos=north west,
]
\addplot[draw=none,fill=shipred,fill opacity=0.15,forget plot] coordinates {
  (0.3496,1.0362) (0.3961,1.1348) (0.4425,1.2271) (0.4890,1.3136) (0.5354,1.3949)
  (0.5819,1.4714) (0.6284,1.5435) (0.6748,1.6117) (0.7213,1.6761) (0.7677,1.7372)
  (0.8142,1.7951) (0.8606,1.8502) (0.9071,1.9026) (0.9535,1.9524) (1.0000,2.0000)
  (1.0000,4.2000) (0.3496,4.2000)
} \closedcycle;
\addplot[seagreen,thick] coordinates {
  (0.0000,0.0000) (0.0699,0.2614) (0.1398,0.4908) (0.2098,0.6936) (0.2797,0.8743)
  (0.3496,1.0362) (0.3961,1.1348) (0.4425,1.2271) (0.4890,1.3136) (0.5354,1.3949)
  (0.5819,1.4714) (0.6284,1.5435) (0.6748,1.6117) (0.7213,1.6761) (0.7677,1.7372)
  (0.8142,1.7951) (0.8606,1.8502) (0.9071,1.9026) (0.9535,1.9524) (1.0000,2.0000)
};
\addlegendentry{downtime floor $W_\infty=\xi/(\eta(\xi+\eta))$}
% Third series: the ACTUAL W_bd at a finite skill premium, so that "infimum,
% never attained" is visible rather than asserted -- the property Theorem 3's
% boxed inequality turns on. mu_s=5 is the paper's own canary instance; recomputed
% here from the box's identity W_bd = 1/(mu_eff-lam) + mu_eff/(mu_eff-lam)*W_inf
% with mu_eff = mu_s*eta/(xi+eta), lam=1, eta=0.25, xi=eta*(xi/eta). Spot check
% against the paper's stated canary: at xi/eta=2 (xi=0.5) this gives 8.1667,
% matching the "truth is 8.17" quoted in Sec. "The succession price".
\addplot[black!45,thick,densely dotted] coordinates {
  (0.0000,0.2500) (0.0699,0.6047) (0.1398,0.9307) (0.2098,1.2343) (0.2797,1.5190)
  (0.3496,1.7890) (0.3961,1.9619) (0.4425,2.1301) (0.4890,2.2948) (0.5354,2.4561)
  (0.5819,2.6152) (0.6284,2.7721) (0.6748,2.9271) (0.7213,3.0812) (0.7677,3.2341)
  (0.8142,3.3869) (0.8606,3.5393) (0.9071,3.6923) (0.9535,3.8455) (1.0000,4.0000)
};
\addlegendentry{actual $W_{\mathrm{bd}}$ at $\mu_s{=}5$ (always above the floor)}
\addplot[harborblue,thick,dashed] coordinates {(0,1.0362) (1.0,1.0362)};
\addlegendentry{pool cost line $K=A/(w\lambda)+W_{\mathrm{pool}}$}
\addplot[black,dotted,forget plot] coordinates {(0.34962,0) (0.34962,4.2)};
\node[font=\scriptsize,rotate=90,anchor=south west] at (axis cs:0.3646,0.25) {$D^\star=0.350$};
% Moved down from y=2.15: legend pos=north west with these two long entries
% (one with a stacked fraction) renders nearly full plot width and reaches
% down to roughly y=2.0-2.1, which the original placement (just under
% ymax=2.6) sat inside -- not a shared-style issue, this callout is local
% to the panel. Lower + smaller also clears the mu_s marker/label block
% (~y=1.12-1.48) below it.
\addplot[shipred,mark=*,mark size=2.2pt,only marks,forget plot] coordinates {(0.5244,1.3761)};
\node[shipred,font=\scriptsize,anchor=north west,align=left,fill=white,fill opacity=0.85,inner sep=1.5pt] at (axis cs:0.56,0.85) {$\mu_s{=}10^8$ still loses\\($W{=}1.376{>}K$)};
\draw[shipred,->,thick] (axis cs:0.56,0.83) -- (axis cs:0.53,1.32);
\end{axis}
\end{tikzpicture}
\end{minipage}
\caption{The regime diagram. Panel A: the exact boundary $g(\rho,2)=1+2\rho-\rho^2$ (solid) against the falsified proposal
$\tilde g$ (dashed). The two hatched lenses --- each labelled on the figure --- are the two error regimes, and the proposal is
wrong on \emph{both} sides of the crossing at $\rho=(3-\sqrt5)/2$, so no constant rescues it; the refuting instances are marked
($r{=}1.15$ at $\rho{=}0.1$; $r{=}1.9$ at $\rho{=}0.5$). The $y$-axis is clipped at $r{=}4$: $\tilde g$ leaves the frame near
$\rho{=}0.76$ and reaches $12.5$ by $\rho{=}0.92$, which is the divergence $g$ does not have. Panel B:
the succession price at the reference instance ($\eta{=}0.25$, $K{=}1.036$): the downtime floor $W_\infty=\xi/(\eta(\xi+\eta))$
crosses the pool cost line at $D^\star=0.350$, and to its right pooling wins at every \emph{finite} skill premium. The dotted
grey series is the actual $W_{\mathrm{bd}}$ at a finite $\mu_s{=}5$, drawn to make the floor's un-attainedness visible: no
specialist ever sits on it, which is why Theorem 3's criterion is non-strict at $D^\star$. The marked point is a
$\mu_s=10^8$ specialist still losing at $\xi/\eta=1.5D^\star$ ($W=1.376>K$). Both panels are closed-form curves plotted from
\texttt{figures/src/paper6\_figures.py}; the underlying boundary and $D^\star$ results are \texttt{b8\_specialization.py}
[internal, seed 20260816].}
\label{fig:regime}
\end{figure}
